%(BEGIN_QUESTION)
% Copyright 2007, Tony R. Kuphaldt, released under the Creative Commons Attribution License (v 1.0)
% This means you may do almost anything with this work of mine, so long as you give me proper credit

En dag på jobb møter du en automatisk regulert prosess som ser ut til å oppføre seg dårlig. Prosessvariabelen (PV) driver betydelig bort fra setpunktet (SP) over tid, slik trendskriveren viser. Etter å ha lært teknikker for PID-innstilling kan første instinkt være å begynne å justere regulatorens P-, I- og D-konstanter. Men du vet bedre enn å kaste deg over tuning av en sløyfe før du forstår mer om den!

\vskip 10pt

Identifiser noen faktorer {\it utover} feilinnstilt regulator som kan være ansvarlige for den dårlige reguleringskvaliteten, og hvordan du kan identifisere dem.

\underbar{file i01645}
%(END_QUESTION)





%(BEGIN_ANSWER)

Det finnes mange mulige kilder til ustabilitet i en reguleringssløyfe, blant annet (men ikke begrenset til):

\begin{itemize}
\item{}Feil i primærsensorelementet (PSE), som sender ustabile eller unøyaktige signaler til transmitteren.
\item{}Feil i transmitteren, som sender regulatoren et ustabilt eller unøyaktig prosessvariabelsignal.
\item{}Feil i sluttregulerende element, som ikke responderer korrekt på kommandoer fra regulatoren.
\item{}Feil i signalvei(er), som ikke overfører signaler korrekt til og fra regulatoren.
\item{}Feilkonstruert sløyfe, der stabil regulering i praksis er {\it umulig} å oppnå slik anlegget er bygget.
\item{}Uvanlige variasjoner i prosesslast, som utfordrer instrumenteringens evne til å måle og regulere.
\end{itemize}

%(END_ANSWER)





%(BEGIN_NOTES)

Denne listen er naturligvis ikke uttømmende. Den viser likevel at man aldri bør anta at regulatoren er kilden til sløyfeustabilitet bare fordi den er den mest tilgjengelige komponenten i sløyfen.

%INDEX% Control, PID tuning: causes of instability other than controller tuning

%(END_NOTES)
